% 2023-04-10 v1.1
% 1. 编辑以下相关文件：
% 按Ctrl键同时用鼠标左键单击以下文件名可打开相关文件
% % 学校代码
\schoolcode{10459}
% 学号或申请号
\id{XXXXXX}
% 密级
\secretlevel{}
% 论文题目
\ctitle{XXXXXXX}
\etitle{XXXXXXX} 
% 学生姓名
\cauthor{XXX} 
\eauthor{XXX}
% 所在院系（只有英文）
\edepartment{School of Materials Science and Engineering}
% 指导老师
\csupervisor{XXX}
\protitle{XX}%导师职称
\esupervisor{XXX} 
% 第二导师
\ccosupervisor{}
\coprotitle{}%导师职称
% 学科门类
\csubject{工学}
% 专业名称
\cmajor{控制科学与工程}
\emajor{Materials Science and Engineering} 
% 完成时间
\cdate{2023年7月}
\edate{July, 2023}

% \begin{cabstract}
应概括地反映出本论文的主要内容，具有独立性和自创性，包括研究工作目的、方法、结果和结论，要突出本论文的创造性成果。摘要力求语言精炼准确，硕士学位论文建议1000字以内，博士学位论文建议2000字以内。部分学生用外文撰写学位论文时，博士学位论文的摘要应不少于5000字符，硕士应不少于2500字符。摘要中不可出现图片、图表、表格或其他插图材料。

关键词是为了便于做文献索引和检索工作而从论文中选取出来用以表示全文主题内容信息的单词或术语。关键词应体现论文特色，具有语义性，在论文中有明确的出处，并应尽量采用《汉语主题词表》或各专业主题词表提供的规范词。在摘要内容下方另起一行标明，一般3～8个，之间用‘；’分开。为便于国际交流，应标注与中文对应的英文关键词。
\end{cabstract}
\ckeywords{关键词1,关键词2,关键词3,关键词4,关键词5}
\begin{eabstract} 
The abstract content corresponds to the Chinese abstract. It should be noted that the case of foreign language and the writing in italics should meet the relevant requirements.

Abstract内容与中文摘要相对应，应注意外文大小写、正斜体书写符合有关要求。
\end{eabstract}

\ekeywords{Keywords1,Keywords2,Keywords3,Keywords4,Keywords5}

% %符号说明样式
% \denotationstyle                
% 符号说明环境
\begin{denotation}{3em}
    \item[$m$] 质量
    \item[$E$] 能量
    \item[$c$] 光速
    \item[$v$] 速度
\end{denotation}

% % chapter1.tex
\chapter{内容要求}
\label{cha:overview}

研究生学位论文使用汉字（除留学生及外国语言文学专业要求用其它文字外）撰写。学位论文一般由十五部分组成，依次为：

（1）封面；

（2）英文封面；

（3）学位论文原创性声明；

（4）学位论文使用授权声明；

（5）摘要；

（6）Abstract；

（7）目录；

（8）图和附表清单；

（9）符号说明；

（10）正文；

（11）注释；

（12）参考文献；

（13）附录（A，B，…）；

（14）个人简历、在学期间发表的学术论文及研究成果；

（15）致谢。


% %chapter2.tex
\chapter{格式要求}
\label{geshiyaoqiu}
学位论文每部分一级标题从新的一页开始，各部分要求如下：

\section{封面}
封面是学位论文的外表面，对论文起装潢和保护作用，并提供相关的信息。博士学位论文、硕士学位论文、专业博士学位论文、专业硕士学位论文使用不同封面。

学校代码：10459

学号或申请号：全日制研究生填写学号，在职申请学位人员填写申请号。

密级：非涉密（公开）论文不需标注密级，涉密论文须标注论文的密级（内部、秘密、或机密），同时还应注明相应的保密年限。

论文题目：应简明扼要地概括和反映出论文的核心内容，一般不宜超过25字，必要时可加副标题，若本论文为国家自然基金等资助项目可在题目下一行用宋体12磅居中注明。

培养院系：填写所属院（系）的全名。

学科门类、专业名称：以国务院学位委员会1997年颁布的《授予博士、硕士学位和培养研究生的学科、专业目录》为准，学科门类以十二大门类为准填写，专业名称以二级学科为准填写。

专业学位名称（仅限于专业博士、硕士学位论文）：填写所申请的专业硕士学位的全称，如工商管理硕士、工程硕士等。

导师姓名：填写导师的姓名、职称（教授、研究员等）

完成时间：填写论文成文打印日期。格式：例2010年5月；2010年12月

\section{英文封面}
按照中文封面格式在封面相应位置填写英文名称，含中文封面的主要内容。

\section{学位论文原创性声明}
本部分内容直接从郑州大学研究生院资料下载中下载，提交时学位论文作者须亲笔签名。

\section{学位论文使用授权声明}
本部分内容直接从郑州大学研究生院资料下载中下载，提交时学位论文作者须亲笔签名。

\section{摘要}
应概括地反映出本论文的主要内容，具有独立性和自创性，包括研究工作目的、方法、结果和结论，要突出本论文的创造性成果。摘要力求语言精炼准确，硕士学位论文建议1000字以内，博士学位论文建议2000字以内。部分学生用外文撰写学位论文时，博士学位论文的摘要应不少于5000字符，硕士应不少于2500字符。摘要中不可出现图片、图表、表格或其他插图材料。

关键词是为了便于做文献索引和检索工作而从论文中选取出来用以表示全文主题内容信息的单词或术语。关键词应体现论文特色，具有语义性，在论文中有明确的出处，并应尽量采用《汉语主题词表》或各专业主题词表提供的规范词。在摘要内容下方另起一行标明，一般3～8个，之间用‘；’分开。为便于国际交流，应标注与中文对应的英文关键词。


\section{Abstract}
Abstract内容与中文摘要相对应，应注意外文大小写、正斜体书写符合有关要求。

\section{目录}
目录应将文内的章节标题依次排列。

\section{图和附表清单}
图表较多时使用。图的清单应有序号、图题和页码。表的清单应有序号、表题和页码。

\section{符号说明}
如果论文中使用了大量的符号、标志、缩略词、首字母缩写、专门计量单位、自定义名词和术语等，应编写成注释说明汇集表。若上述符号使用数量不多，可以不设此部分，但必须在论文中出现时加以说明。

\section{正文}
正文是学位论文的主体和核心部分，每一章一级标题应另起页，一般包括以下几个方面：

（1）引言（第1章）：包括研究的目的和意义，问题的提出，选题的背景，文献综述，研究方法，论文结构安排等。

（2）各具体章节：本部分是论文作者的研究内容，是论文的核心。各章之间互相关联，符合逻辑顺序。

（3）结论（最后1章）：是学位论文最终和总体的结论，不是正文中各段小结的简单重复，应明确、精练、完整、准确。着重阐述作者的创造性工作及所取得的研究成果在本学术领域的地位、作用和意义，结论应包括论文的核心观点，交代研究工作的局限，提出未来工作的意见或建议。


\section{注释}
当论文中的字、词或短语，需要进一步加以说明，而又没有具体的文献来源时，用注释。注释一般在社会科学中用得较多。应控制论文中的注释数量，不宜过多。由于论文篇幅较长，建议采用文中编号加“脚注”的方式。

\section{参考文献}
参考文献表是文中引用的有具体文字来源的文献集合。为了反映论文的科学依据和作者尊重他人研究成果的严肃态度以及向读者提供有关信息的出处，应列出参考文献表。参考文献表中列出的一般应限于作者直接阅读过的、最主要的、发表在正式出版物上的文献。私人通信和未公开发表的资料，一般不宜列入参考文献，可紧跟在引用的内容之后在文中编号，用“脚注”的方式标注在当页的下方。正文中引用参考文献的部位，须用上标标注[参考文献序号]，如 \cite{chen2001,nadkarni1992,hua1973,zhu1973,huo1981,timoshenko1959,zhang1991,dupont1974,ding2001,cairns1965,patent,standard,Karl2022}。参考文献应按文中引用出现的顺序排列。

\section{附录}
有些材料编入文章主体会有损于编排的条理性和逻辑性，或有碍于文章结构的紧凑和突出主题思想等，可将这些材料作为附录编排于全文的末尾。

附录的序号用A，B，C…系列，如附录A，附录B…。附录中的公式、图和表的编号分别用A1，A2…系列；图A1，图A2…系列；表A1，表A2…系列。每个附录应有标题。


附录的序号用A，B，C…系列，如附录系列，如附录A，附录B…。附录中的公式、图。附录中的公式、图和表的编号分别用和表的编号分别用A1，A2…系列；图A1，图A2…系列；表系列；表A1，表A2…系列。每个附录应有标题。每个附录应有标题。

\section{个人简历、在学期间发表的学术论文与研究成果}
个人简历包括出生年月日、获得学士、硕士学位的学校、时间等；学术论文研究成果按发表的时间顺序列出；研究成果可以是在学期间参加的研究项目、申请的专利或获奖情况等。

\section{致谢}
致谢是作者对该论文的形成作过贡献的组织或个人及参考文献的作者予以感谢的文字记载，语言要诚恳、恰当、简短。


% % chpater3.tex

\chapter{书写要求}
\label{shuxieyaoqiu}
\section{论文字数、文字、标点符号和数字}

硕士学位论文要求一般不少于3万字（专业学位论文等特殊类别按相关规定执行），博士学位论文\footnote{脚注示例。}要求不少于5万字。

除留学生及外语专业研究生外，学位论文一律用汉字书写。汉字的使用应严格执行国家的有关规定，除特殊需要外，不得使用已废除的繁体字、异体字等不规范汉字。标点符号的用法应该以GB/T 15834—1995《标点符号用法》为准。数字用法应该以GB/T 15835—1995《出版物上数字用法的规定》为准，单位用法以GB 3100～3102—93《国际单位制及其应用》为准。

留学生的学位论文所采用语种可以和导师商定，但论文封面须用中文。

\section{密级}

各密级的最长保密年限及书写格式规定如下：

（1）内部。5年（最长5年，可少于5年）。

（2）秘密\lower .2ex\hbox {\FiveStar}。10年（最长10年，可少于10年）。

（3）机密\lower .2ex\hbox {\FiveStar}。20年（最长20年，可少于20年）。

\section{层次标题}	

层次标题要简短明确，同一层次的标题应尽可能“排比”，词（或词组）类型相同（或相近），意义相关，语气一致。 

各层次标题一律用阿拉伯数字连续编号；不同层次的数字之间用英文输入状态下小圆点“．”相隔，末位数字后面不加点号，如“1”，“2.1”，“3.1.2”等；各层次的序号均左顶格起排，编号与标题或文字间空一个汉字的间隙。段的文字空两个汉字起排，回行时顶格排。如：

1\hspace{\ccwd}$\times\times\times\times$（一级标题）

1.1\hspace{\ccwd}$\times\times\times\times$（二级标题）

1.1.1\hspace{\ccwd}$\times\times\times\times$（三级标题）

\section{篇眉和页码}

篇眉从中文摘要开始，内容与该部分的一级标题相同。

页码从正文开始按阿拉伯数字（1，2，3……）连续编排，之前的部分(中文摘要、Abstract、目录等)用大写罗马数字（Ⅰ，Ⅱ，Ⅲ……）单独编排，均居中排列。

\section{有关图、表、表达式}

\subsection{图}

图包括曲线图、构造图、示意图、框图、流程图、记录图、地图、照片等。

图要精选，要具有自明性，切忌与表及文字表述重复。

图要清楚，但坐标比例不要过分放大，同一图上不同曲线的点要分别用不同形状的标识符标出。图中的术语、符号、单位等应与正文表述中所用一致。图在文中的布局要合理，一般随文编排，先见文字后见图。

图序与图题：图序一律采用阿拉伯数字分章编号，第3章第1个图的图序为“\reffig{fig3_1}”；图题应简明。图序和图题间空1个汉字距，居中排于图的下方。例如：

\begin{figure}[h]
  \centering
  \includegraphics{示例图1}
  \caption{非线性构形状态转移过程示意图}
  \label{fig3_1}
\end{figure}

照片图要求主题和主要显示部分的轮廓鲜朗，便于制版，要标明出处。如用放大缩小的复制品，必须清晰，反差适中。照片上应有表示目的物尺寸的标度。

\subsection{表}

表中参数应标明量和单位的符号。表一般随文排，先见相应文字后见表。

表序与表题：表序一律采用阿拉伯数字分章编号，如第3章第1个表的表序表示为“\reftab{tab3_1}”；表题应简明。表序和表题间空1个汉字距，居中排于表的上方。

表的编排，一般是内容和测试项目由左至右横读，数据依序竖读。表的编排建议采用国际通行的三线表。例如：           

\begin{table}[h]
\caption{文献类型和标识代码}  
\label{tab3_1}
\centering
\tabcolsep=20pt
\begin{tabular}{cccc}
\toprule
文献类型 & 标识代码 & 文献类型 & 标识代码\\
\midrule
普通图书 & M & 会议录 & C\\
汇编 & G & 报纸 & N\\
期刊 & J & 学位论文 &     D\\
报告 & R & 标准 & S\\
专利 & P & 数据库 & DB\\
计算机程序 & CP & 电子公告 & EB\\
\bottomrule
\end{tabular}
\end{table}

表格较大，不能在一页打印、需要转页排时，需在续表上方居中注明“续表”，续表的表头应重复排出。

\subsection{表达式}

表达式主要指数学表达式，也包括文字表达式。表达式需另行起排，并用阿拉伯数字分章编号。序号加圆括号，右顶格排。例如：第3章第1个表达式为\refequ{equ3_1}。
\begin{equation}
  \label{equ3_1}
  \Phi(\Theta )=
  \begin{cases}
    0& \text{if $E=E_0$},\\
    1& \text{if $E=E_1$}.
  \end{cases}
\end{equation}

\section{参考文献}

文后著录的参考文献务必实事求是。论文中引用过的文献必须著录，未引用的文献不得出现。应遵循学术道德规范，避免涉嫌抄袭、剽窃等学术不端行为。

参考文献一般应是作者亲自考察过的对学位论文有参考价值的文献，除特殊情况外，一般不应间接引用。

文献的作者不超过3位时，全部列出；超过3位时，只列前3位，后面加“等”字或相应的外文；作者姓名之间用“，”分开。建议根据《中国高校自然科学学报编排规范》的要求书写参考文献，并按顺序编码制，即按中文引用的参考顺序将参考文献附于文末。

几种主要参考文献著录表的格式为：

连续出版物: [序号] 作者.文题[文献类型标识].刊名，年，卷号（期号）：起～止页码

专（译）著:  [序号] 作者.书名（，译者）[文献类型标识].出版地：出版者，出版年.起～止页码

论 文 集: [序号] 作者.文题[文献类型标识].见（in）：编者，编（eds）.文集名.出版

地：出版者，出版年.起～止页码

学 位 论 文：[序号] 姓名.文题[文献类型标识].[XX学位论文].授予单位所在地：授予单

位，授予年

专      利：[序号] 申请者.专利名[文献类型标识].国名，专利文献种类，专利号，授权

公告日

技 术 标 准：[序号] 发布单位.技术标准代号.技术标准名称[文献类型标识].出版地：出版者，出版日期

电子文献：作者.文题[文献类型标识].出版年（更新和修改日期）.获取和访问路径

例如：

[1] 陈建军，车建文，陈勇.具有频率和振型概率约束的工程结构动力优化设计[J].计算力学学报，2001，18（1）:74～80

[2] Nadkarni M A，Nair C K K，Pandey V N，et al．Characterzation of alpha-galactosidase from corynebacterium murisepticum and mechanism of its induction[J].J Gen App Microbiol，1992，38：223～234 

[3] 华罗庚，王 元．论一致分布与近似分析：数论方法（Ⅰ）[J].中国科学，1973,（4）：339～357 

[4] 竺可桢．物候学[M].北京：科学出版社，1973.104～107

[5] 霍夫斯塔主编．禽病学：下册．第7版[M]．胡祥壁译．北京：农业出版社，1981．798～799 

[6] Timoshenko S P．Theory of plate and shells．2nd ed[M]．New York：McGraw-Hil1，1959．17～36

[7] 张全福，王里青．“百家争鸣”与理工科学报编辑工作[G]．见：郑福寿主编．学报编辑论丛：第2集．南京：河海大学出版社，1991：1～４ 

[8] Dupont B．Bone marrow transplantation in severe combined immunodeficiency with an unrclated MLC compatible donor[C]．In：White H J，Smith R，eds．Proceedings of the Third Annual Meeting of the International Society for Experimental Hematology．Houston：International Society for Experimental Hematology，1974．44～46 

[9] 丁光莹.钢筋混凝土框架结构非线性反应分析的随机模拟分析[D].[博士学位论文].上海：同济大学，2001 

[10] Cairns R B．Infrared spectroscopic studies on solid oxygen[D].[dissertation]．Berkeley：Univ of California，1965 

[11] 姜锡洲．一种温热外敷药制备方法[P].中国专利，881056073．1989-07-26 

[12] 全国文献工作标准化技术委员会第六分委员会．GB 6447—86 文摘编写规则[S].北京：中国标准出版社，1986 

[13] 王文惠，孟兵等.利用不变量进行基于内容的图像检索[EB].2005.

http://www.image2003.com/paper/down/2005528115935222.pdf

\section{量和单位}

要严格执行GB 3100～3102—93(国家技术监督局1993-12-27发布，1994-07-01实施)有关量和单位的规定。

量的符号一般为单个拉丁字母或希腊字母，并一律采用斜体（pH例外）。为区别不同情况，可在量符号上附加角标。如：

$A$	--截面积,散热面积；

$T_h$	--弧柱温度；

$B$	--磁感应强度；

$t$	--时间；

$B_r$	--剩磁感应强度；

$t_c$	--触动时间；

$B_s$	--饱和磁感应强度；

$t-d$	--运动时间；

$C$	--电容；

$U$,$u$	--电压。

在表达量值时，在公式、图、表和文字叙述中，一律使用单位的国际符号，且无例外地用正体。单位符号与数值间要留适当间隙。


\section{印刷及装订要求}
论文自中文摘要起正文部分双面印刷，之前部分单面印刷。论文必须用线装或热胶装订，不能使用钉子装订。
% \input{../Common/Demo/模板使用说明}
% \input{../Common/BibFile/BibFile.bib}
% \begin{appendix}

  \chapter{附录说明}

  有些材料编入文章主体会有损于编排的条理性和逻辑性，或有碍于文章结构的紧凑和突出主题思想等，可将这些材料作为附录编排于全文的末尾。

  附录的序号用A B C… 系列，如附录 A ，附录 B… 。附录中的公式、图和表的编号分别用 A1 A2… 系列；图 A1 ，图 A2… 系列；表 A1 ，表 A2… 系列。每个附录应有标题。

  \chapter{第二个附录}

  第二个附录。依次类推。

\end{appendix}
% app01.tex

% \input{../Common/TexFiles/作者简介}
% \begin{ack}
致谢是作者对该论文的形成作过贡献的组织或个人及参考文献的作者予以感谢的文字记载，语言要诚恳、恰当、简短。
\end{ack}


% 2. 编译环境：
% （VScode(加上LaTex Workshop扩展)+Texlive环境）

% 3. 编译生成PDF文件方法：
% （1）在LaTex扩展中“Build LaTex project”中点击“Recipe:xelatex”
% （2）或在些文档的右上角点击绿色“三角播放”按钮。
%  以上两种方法适合快速编译，不能正确显示各类引用编号。
% （3）在LaTex扩展中“Build LaTex project”中点击“Recipe:All”
%  以上方法生成最终包含正确引用的PDF文件。

% 使用“betterformat”选项
%  (1)封面上使用楷体
%  (2)封面“硕士学位论文”和“郑州大学”手写体严格居中（建议选择）
% 不用“betterformat”选项
%  (1)封面使用GB_2312楷体, Windows 10中需要自行下载安装
%  (2)封面“硕士学位论文”和“郑州大学”手写体整体右移0.43cm
% “单面打印”选项
%  按照学校模板要求，硕士论文可以单面打印，需选择“单面打印”
%  如果页码较多，可以去掉该选项，使用双面打印模式。
\documentclass[betterformat,单面打印]{../Common/Template/MasterThesis}

\begin{document}
% 学校代码
\schoolcode{10459}
% 学号或申请号
\id{XXXXXX}
% 密级
\secretlevel{}
% 论文题目
\ctitle{XXXXXXX}
\etitle{XXXXXXX} 
% 学生姓名
\cauthor{XXX} 
\eauthor{XXX}
% 所在院系（只有英文）
\edepartment{School of Materials Science and Engineering}
% 指导老师
\csupervisor{XXX}
\protitle{XX}%导师职称
\esupervisor{XXX} 
% 第二导师
\ccosupervisor{}
\coprotitle{}%导师职称
% 学科门类
\csubject{工学}
% 专业名称
\cmajor{控制科学与工程}
\emajor{Materials Science and Engineering} 
% 完成时间
\cdate{2023年7月}
\edate{July, 2023}

\begin{cabstract}
应概括地反映出本论文的主要内容，具有独立性和自创性，包括研究工作目的、方法、结果和结论，要突出本论文的创造性成果。摘要力求语言精炼准确，硕士学位论文建议1000字以内，博士学位论文建议2000字以内。部分学生用外文撰写学位论文时，博士学位论文的摘要应不少于5000字符，硕士应不少于2500字符。摘要中不可出现图片、图表、表格或其他插图材料。

关键词是为了便于做文献索引和检索工作而从论文中选取出来用以表示全文主题内容信息的单词或术语。关键词应体现论文特色，具有语义性，在论文中有明确的出处，并应尽量采用《汉语主题词表》或各专业主题词表提供的规范词。在摘要内容下方另起一行标明，一般3～8个，之间用‘；’分开。为便于国际交流，应标注与中文对应的英文关键词。
\end{cabstract}
\ckeywords{关键词1,关键词2,关键词3,关键词4,关键词5}
\begin{eabstract} 
The abstract content corresponds to the Chinese abstract. It should be noted that the case of foreign language and the writing in italics should meet the relevant requirements.

Abstract内容与中文摘要相对应，应注意外文大小写、正斜体书写符合有关要求。
\end{eabstract}

\ekeywords{Keywords1,Keywords2,Keywords3,Keywords4,Keywords5}

\frontmatter
%符号说明样式
% \denotationstyle                
% 符号说明环境
\begin{denotation}{3em}
    \item[$m$] 质量
    \item[$E$] 能量
    \item[$c$] 光速
    \item[$v$] 速度
\end{denotation}

\mainmatter
% chapter1.tex
\chapter{内容要求}
\label{cha:overview}

研究生学位论文使用汉字（除留学生及外国语言文学专业要求用其它文字外）撰写。学位论文一般由十五部分组成，依次为：

（1）封面；

（2）英文封面；

（3）学位论文原创性声明；

（4）学位论文使用授权声明；

（5）摘要；

（6）Abstract；

（7）目录；

（8）图和附表清单；

（9）符号说明；

（10）正文；

（11）注释；

（12）参考文献；

（13）附录（A，B，…）；

（14）个人简历、在学期间发表的学术论文及研究成果；

（15）致谢。


%chapter2.tex
\chapter{格式要求}
\label{geshiyaoqiu}
学位论文每部分一级标题从新的一页开始，各部分要求如下：

\section{封面}
封面是学位论文的外表面，对论文起装潢和保护作用，并提供相关的信息。博士学位论文、硕士学位论文、专业博士学位论文、专业硕士学位论文使用不同封面。

学校代码：10459

学号或申请号：全日制研究生填写学号，在职申请学位人员填写申请号。

密级：非涉密（公开）论文不需标注密级，涉密论文须标注论文的密级（内部、秘密、或机密），同时还应注明相应的保密年限。

论文题目：应简明扼要地概括和反映出论文的核心内容，一般不宜超过25字，必要时可加副标题，若本论文为国家自然基金等资助项目可在题目下一行用宋体12磅居中注明。

培养院系：填写所属院（系）的全名。

学科门类、专业名称：以国务院学位委员会1997年颁布的《授予博士、硕士学位和培养研究生的学科、专业目录》为准，学科门类以十二大门类为准填写，专业名称以二级学科为准填写。

专业学位名称（仅限于专业博士、硕士学位论文）：填写所申请的专业硕士学位的全称，如工商管理硕士、工程硕士等。

导师姓名：填写导师的姓名、职称（教授、研究员等）

完成时间：填写论文成文打印日期。格式：例2010年5月；2010年12月

\section{英文封面}
按照中文封面格式在封面相应位置填写英文名称，含中文封面的主要内容。

\section{学位论文原创性声明}
本部分内容直接从郑州大学研究生院资料下载中下载，提交时学位论文作者须亲笔签名。

\section{学位论文使用授权声明}
本部分内容直接从郑州大学研究生院资料下载中下载，提交时学位论文作者须亲笔签名。

\section{摘要}
应概括地反映出本论文的主要内容，具有独立性和自创性，包括研究工作目的、方法、结果和结论，要突出本论文的创造性成果。摘要力求语言精炼准确，硕士学位论文建议1000字以内，博士学位论文建议2000字以内。部分学生用外文撰写学位论文时，博士学位论文的摘要应不少于5000字符，硕士应不少于2500字符。摘要中不可出现图片、图表、表格或其他插图材料。

关键词是为了便于做文献索引和检索工作而从论文中选取出来用以表示全文主题内容信息的单词或术语。关键词应体现论文特色，具有语义性，在论文中有明确的出处，并应尽量采用《汉语主题词表》或各专业主题词表提供的规范词。在摘要内容下方另起一行标明，一般3～8个，之间用‘；’分开。为便于国际交流，应标注与中文对应的英文关键词。


\section{Abstract}
Abstract内容与中文摘要相对应，应注意外文大小写、正斜体书写符合有关要求。

\section{目录}
目录应将文内的章节标题依次排列。

\section{图和附表清单}
图表较多时使用。图的清单应有序号、图题和页码。表的清单应有序号、表题和页码。

\section{符号说明}
如果论文中使用了大量的符号、标志、缩略词、首字母缩写、专门计量单位、自定义名词和术语等，应编写成注释说明汇集表。若上述符号使用数量不多，可以不设此部分，但必须在论文中出现时加以说明。

\section{正文}
正文是学位论文的主体和核心部分，每一章一级标题应另起页，一般包括以下几个方面：

（1）引言（第1章）：包括研究的目的和意义，问题的提出，选题的背景，文献综述，研究方法，论文结构安排等。

（2）各具体章节：本部分是论文作者的研究内容，是论文的核心。各章之间互相关联，符合逻辑顺序。

（3）结论（最后1章）：是学位论文最终和总体的结论，不是正文中各段小结的简单重复，应明确、精练、完整、准确。着重阐述作者的创造性工作及所取得的研究成果在本学术领域的地位、作用和意义，结论应包括论文的核心观点，交代研究工作的局限，提出未来工作的意见或建议。


\section{注释}
当论文中的字、词或短语，需要进一步加以说明，而又没有具体的文献来源时，用注释。注释一般在社会科学中用得较多。应控制论文中的注释数量，不宜过多。由于论文篇幅较长，建议采用文中编号加“脚注”的方式。

\section{参考文献}
参考文献表是文中引用的有具体文字来源的文献集合。为了反映论文的科学依据和作者尊重他人研究成果的严肃态度以及向读者提供有关信息的出处，应列出参考文献表。参考文献表中列出的一般应限于作者直接阅读过的、最主要的、发表在正式出版物上的文献。私人通信和未公开发表的资料，一般不宜列入参考文献，可紧跟在引用的内容之后在文中编号，用“脚注”的方式标注在当页的下方。正文中引用参考文献的部位，须用上标标注[参考文献序号]，如 \cite{chen2001,nadkarni1992,hua1973,zhu1973,huo1981,timoshenko1959,zhang1991,dupont1974,ding2001,cairns1965,patent,standard,Karl2022}。参考文献应按文中引用出现的顺序排列。

\section{附录}
有些材料编入文章主体会有损于编排的条理性和逻辑性，或有碍于文章结构的紧凑和突出主题思想等，可将这些材料作为附录编排于全文的末尾。

附录的序号用A，B，C…系列，如附录A，附录B…。附录中的公式、图和表的编号分别用A1，A2…系列；图A1，图A2…系列；表A1，表A2…系列。每个附录应有标题。


附录的序号用A，B，C…系列，如附录系列，如附录A，附录B…。附录中的公式、图。附录中的公式、图和表的编号分别用和表的编号分别用A1，A2…系列；图A1，图A2…系列；表系列；表A1，表A2…系列。每个附录应有标题。每个附录应有标题。

\section{个人简历、在学期间发表的学术论文与研究成果}
个人简历包括出生年月日、获得学士、硕士学位的学校、时间等；学术论文研究成果按发表的时间顺序列出；研究成果可以是在学期间参加的研究项目、申请的专利或获奖情况等。

\section{致谢}
致谢是作者对该论文的形成作过贡献的组织或个人及参考文献的作者予以感谢的文字记载，语言要诚恳、恰当、简短。


% chpater3.tex

\chapter{书写要求}
\label{shuxieyaoqiu}
\section{论文字数、文字、标点符号和数字}

硕士学位论文要求一般不少于3万字（专业学位论文等特殊类别按相关规定执行），博士学位论文\footnote{脚注示例。}要求不少于5万字。

除留学生及外语专业研究生外，学位论文一律用汉字书写。汉字的使用应严格执行国家的有关规定，除特殊需要外，不得使用已废除的繁体字、异体字等不规范汉字。标点符号的用法应该以GB/T 15834—1995《标点符号用法》为准。数字用法应该以GB/T 15835—1995《出版物上数字用法的规定》为准，单位用法以GB 3100～3102—93《国际单位制及其应用》为准。

留学生的学位论文所采用语种可以和导师商定，但论文封面须用中文。

\section{密级}

各密级的最长保密年限及书写格式规定如下：

（1）内部。5年（最长5年，可少于5年）。

（2）秘密\lower .2ex\hbox {\FiveStar}。10年（最长10年，可少于10年）。

（3）机密\lower .2ex\hbox {\FiveStar}。20年（最长20年，可少于20年）。

\section{层次标题}	

层次标题要简短明确，同一层次的标题应尽可能“排比”，词（或词组）类型相同（或相近），意义相关，语气一致。 

各层次标题一律用阿拉伯数字连续编号；不同层次的数字之间用英文输入状态下小圆点“．”相隔，末位数字后面不加点号，如“1”，“2.1”，“3.1.2”等；各层次的序号均左顶格起排，编号与标题或文字间空一个汉字的间隙。段的文字空两个汉字起排，回行时顶格排。如：

1\hspace{\ccwd}$\times\times\times\times$（一级标题）

1.1\hspace{\ccwd}$\times\times\times\times$（二级标题）

1.1.1\hspace{\ccwd}$\times\times\times\times$（三级标题）

\section{篇眉和页码}

篇眉从中文摘要开始，内容与该部分的一级标题相同。

页码从正文开始按阿拉伯数字（1，2，3……）连续编排，之前的部分(中文摘要、Abstract、目录等)用大写罗马数字（Ⅰ，Ⅱ，Ⅲ……）单独编排，均居中排列。

\section{有关图、表、表达式}

\subsection{图}

图包括曲线图、构造图、示意图、框图、流程图、记录图、地图、照片等。

图要精选，要具有自明性，切忌与表及文字表述重复。

图要清楚，但坐标比例不要过分放大，同一图上不同曲线的点要分别用不同形状的标识符标出。图中的术语、符号、单位等应与正文表述中所用一致。图在文中的布局要合理，一般随文编排，先见文字后见图。

图序与图题：图序一律采用阿拉伯数字分章编号，第3章第1个图的图序为“\reffig{fig3_1}”；图题应简明。图序和图题间空1个汉字距，居中排于图的下方。例如：

\begin{figure}[h]
  \centering
  \includegraphics{示例图1}
  \caption{非线性构形状态转移过程示意图}
  \label{fig3_1}
\end{figure}

照片图要求主题和主要显示部分的轮廓鲜朗，便于制版，要标明出处。如用放大缩小的复制品，必须清晰，反差适中。照片上应有表示目的物尺寸的标度。

\subsection{表}

表中参数应标明量和单位的符号。表一般随文排，先见相应文字后见表。

表序与表题：表序一律采用阿拉伯数字分章编号，如第3章第1个表的表序表示为“\reftab{tab3_1}”；表题应简明。表序和表题间空1个汉字距，居中排于表的上方。

表的编排，一般是内容和测试项目由左至右横读，数据依序竖读。表的编排建议采用国际通行的三线表。例如：           

\begin{table}[h]
\caption{文献类型和标识代码}  
\label{tab3_1}
\centering
\tabcolsep=20pt
\begin{tabular}{cccc}
\toprule
文献类型 & 标识代码 & 文献类型 & 标识代码\\
\midrule
普通图书 & M & 会议录 & C\\
汇编 & G & 报纸 & N\\
期刊 & J & 学位论文 &     D\\
报告 & R & 标准 & S\\
专利 & P & 数据库 & DB\\
计算机程序 & CP & 电子公告 & EB\\
\bottomrule
\end{tabular}
\end{table}

表格较大，不能在一页打印、需要转页排时，需在续表上方居中注明“续表”，续表的表头应重复排出。

\subsection{表达式}

表达式主要指数学表达式，也包括文字表达式。表达式需另行起排，并用阿拉伯数字分章编号。序号加圆括号，右顶格排。例如：第3章第1个表达式为\refequ{equ3_1}。
\begin{equation}
  \label{equ3_1}
  \Phi(\Theta )=
  \begin{cases}
    0& \text{if $E=E_0$},\\
    1& \text{if $E=E_1$}.
  \end{cases}
\end{equation}

\section{参考文献}

文后著录的参考文献务必实事求是。论文中引用过的文献必须著录，未引用的文献不得出现。应遵循学术道德规范，避免涉嫌抄袭、剽窃等学术不端行为。

参考文献一般应是作者亲自考察过的对学位论文有参考价值的文献，除特殊情况外，一般不应间接引用。

文献的作者不超过3位时，全部列出；超过3位时，只列前3位，后面加“等”字或相应的外文；作者姓名之间用“，”分开。建议根据《中国高校自然科学学报编排规范》的要求书写参考文献，并按顺序编码制，即按中文引用的参考顺序将参考文献附于文末。

几种主要参考文献著录表的格式为：

连续出版物: [序号] 作者.文题[文献类型标识].刊名，年，卷号（期号）：起～止页码

专（译）著:  [序号] 作者.书名（，译者）[文献类型标识].出版地：出版者，出版年.起～止页码

论 文 集: [序号] 作者.文题[文献类型标识].见（in）：编者，编（eds）.文集名.出版

地：出版者，出版年.起～止页码

学 位 论 文：[序号] 姓名.文题[文献类型标识].[XX学位论文].授予单位所在地：授予单

位，授予年

专      利：[序号] 申请者.专利名[文献类型标识].国名，专利文献种类，专利号，授权

公告日

技 术 标 准：[序号] 发布单位.技术标准代号.技术标准名称[文献类型标识].出版地：出版者，出版日期

电子文献：作者.文题[文献类型标识].出版年（更新和修改日期）.获取和访问路径

例如：

[1] 陈建军，车建文，陈勇.具有频率和振型概率约束的工程结构动力优化设计[J].计算力学学报，2001，18（1）:74～80

[2] Nadkarni M A，Nair C K K，Pandey V N，et al．Characterzation of alpha-galactosidase from corynebacterium murisepticum and mechanism of its induction[J].J Gen App Microbiol，1992，38：223～234 

[3] 华罗庚，王 元．论一致分布与近似分析：数论方法（Ⅰ）[J].中国科学，1973,（4）：339～357 

[4] 竺可桢．物候学[M].北京：科学出版社，1973.104～107

[5] 霍夫斯塔主编．禽病学：下册．第7版[M]．胡祥壁译．北京：农业出版社，1981．798～799 

[6] Timoshenko S P．Theory of plate and shells．2nd ed[M]．New York：McGraw-Hil1，1959．17～36

[7] 张全福，王里青．“百家争鸣”与理工科学报编辑工作[G]．见：郑福寿主编．学报编辑论丛：第2集．南京：河海大学出版社，1991：1～４ 

[8] Dupont B．Bone marrow transplantation in severe combined immunodeficiency with an unrclated MLC compatible donor[C]．In：White H J，Smith R，eds．Proceedings of the Third Annual Meeting of the International Society for Experimental Hematology．Houston：International Society for Experimental Hematology，1974．44～46 

[9] 丁光莹.钢筋混凝土框架结构非线性反应分析的随机模拟分析[D].[博士学位论文].上海：同济大学，2001 

[10] Cairns R B．Infrared spectroscopic studies on solid oxygen[D].[dissertation]．Berkeley：Univ of California，1965 

[11] 姜锡洲．一种温热外敷药制备方法[P].中国专利，881056073．1989-07-26 

[12] 全国文献工作标准化技术委员会第六分委员会．GB 6447—86 文摘编写规则[S].北京：中国标准出版社，1986 

[13] 王文惠，孟兵等.利用不变量进行基于内容的图像检索[EB].2005.

http://www.image2003.com/paper/down/2005528115935222.pdf

\section{量和单位}

要严格执行GB 3100～3102—93(国家技术监督局1993-12-27发布，1994-07-01实施)有关量和单位的规定。

量的符号一般为单个拉丁字母或希腊字母，并一律采用斜体（pH例外）。为区别不同情况，可在量符号上附加角标。如：

$A$	--截面积,散热面积；

$T_h$	--弧柱温度；

$B$	--磁感应强度；

$t$	--时间；

$B_r$	--剩磁感应强度；

$t_c$	--触动时间；

$B_s$	--饱和磁感应强度；

$t-d$	--运动时间；

$C$	--电容；

$U$,$u$	--电压。

在表达量值时，在公式、图、表和文字叙述中，一律使用单位的国际符号，且无例外地用正体。单位符号与数值间要留适当间隙。


\section{印刷及装订要求}
论文自中文摘要起正文部分双面印刷，之前部分单面印刷。论文必须用线装或热胶装订，不能使用钉子装订。
\input{模板使用说明}
\backmatter
\bibliography{../Common/BibFile/BibFile.bib}
\begin{appendix}

  \chapter{附录说明}

  有些材料编入文章主体会有损于编排的条理性和逻辑性，或有碍于文章结构的紧凑和突出主题思想等，可将这些材料作为附录编排于全文的末尾。

  附录的序号用A B C… 系列，如附录 A ，附录 B… 。附录中的公式、图和表的编号分别用 A1 A2… 系列；图 A1 ，图 A2… 系列；表 A1 ，表 A2… 系列。每个附录应有标题。

  \chapter{第二个附录}

  第二个附录。依次类推。

\end{appendix}
% app01.tex

\input{作者简介}
\begin{ack}
致谢是作者对该论文的形成作过贡献的组织或个人及参考文献的作者予以感谢的文字记载，语言要诚恳、恰当、简短。
\end{ack}

\end{document}